% Preview source code for paragraph 13

\begin{table}[!h]
\caption{Fe 在不同状态下的表现形式}

\resizebox{\linewidth}{!}{%

\begin{tabular}{ccccccc}
\toprule 
\textbf{功能} & \textbf{角色} & \textbf{状态} & \textbf{涉及类型} & \textbf{触发器} & \textbf{关键字/核心台词} & \textbf{表现形式 ｜ 例子}\tabularnewline
\midrule
\midrule 
\multirow{13}{*}{Fe} & Hero/Parent & 正常 & - & N/A & 和谐，确认，共情 & 情绪传染，维护关系，社交礼仪\tabularnewline
\cmidrule{2-7}
 & \multirow{2}{*}{Child} & 游戏 & \multirow{2}{*}{ENTP, ESTP} & N/A & \multirow{2}{*}{大家都爱我} & 迷人捣蛋，幽默风趣，活跃气氛\tabularnewline
\cmidrule{3-3}\cmidrule{5-5}\cmidrule{7-7}
 &  & 幼稚 &  & 压力 &  & 渴望关注，哗众取宠，惧怕被弃\tabularnewline
\cmidrule{2-7}
 & \multirow{3}{*}{Anima} & \begin{cellvarwidth}[t]
\centering
正常\\
(Aspiration)
\end{cellvarwidth} & \multirow{3}{*}{INTP, ISTP} & N/A & ``我想被人喜欢'' & \begin{cellvarwidth}[t]
\centering
憨厚微笑，笨拙示好\\
极其客气，送直男礼物，默默帮忙
\end{cellvarwidth}\tabularnewline
\cmidrule{3-3}\cmidrule{5-7}
 &  & \begin{cellvarwidth}[t]
\centering
压抑\\
(Suppression)
\end{cellvarwidth} &  & 防御性屏蔽 & ``社交是Ti的消耗'' & \begin{cellvarwidth}[t]
\centering
社交隐身，情感断连\\
无论谁发信都不回，完全切断人际信号
\end{cellvarwidth}\tabularnewline
\cmidrule{3-3}\cmidrule{5-7}
 &  & \begin{cellvarwidth}[t]
\centering
劫持/抓取\\
(Grip)
\end{cellvarwidth} &  & \begin{cellvarwidth}[t]
\centering
逻辑失效，\\
被群体排斥
\end{cellvarwidth} & \begin{cellvarwidth}[t]
\centering
``没人爱我，\\
我是多余的''
\end{cellvarwidth} & \begin{cellvarwidth}[t]
\centering
情绪崩溃大哭，极其敏感，\\
像个受伤的孩子
\end{cellvarwidth}\tabularnewline
\cmidrule{2-7}
 & \multirow{2}{*}{Nemesis} & \begin{cellvarwidth}[t]
\centering
防御表现\\
(Stubbornness)
\end{cellvarwidth} & \multirow{2}{*}{INFP, ISFP} & \begin{cellvarwidth}[t]
\centering
主导功能\\
受到挑战
\end{cellvarwidth} & \multirow{2}{*}{``别以此绑架我''} & \begin{cellvarwidth}[t]
\centering
拒绝从众：\\
极度反感社交压力和``必须做的事''，\\
故意表现得冷漠
\end{cellvarwidth}\tabularnewline
\cmidrule{3-3}\cmidrule{5-5}\cmidrule{7-7}
 &  & \begin{cellvarwidth}[t]
\centering
投射怀疑\\
(Paranoia)
\end{cellvarwidth} &  & \begin{cellvarwidth}[t]
\centering
对他人该功能\\
莫名厌烦
\end{cellvarwidth} &  & \begin{cellvarwidth}[t]
\centering
怀疑虚伪：\\
认为别人的客套和热情都是有目的的操纵
\end{cellvarwidth}\tabularnewline
\cmidrule{2-7}
 & \multirow{2}{*}{Critic} & 对外打压 & \multirow{2}{*}{ENFP, ESFP} & 防御机制 & \multirow{2}{*}{批判社交: 没人理你} & \begin{cellvarwidth}[t]
\centering
社交孤立，\\
攻击对方情商或者礼仪
\end{cellvarwidth}\tabularnewline
\cmidrule{3-3}\cmidrule{5-5}\cmidrule{7-7}
 &  & 对内自毁 &  & 防御机制 &  & 归属崩塌，总觉自己格格不入\tabularnewline
\cmidrule{2-7}
 & Trickster & 忽视/贬低 & INTJ, ISTJ & 无意识的黑洞 & \begin{cellvarwidth}[t]
\centering
``虚伪的客套\\
有什么用？\\
讲重点。''
\end{cellvarwidth} & \begin{cellvarwidth}[t]
\centering
视社交礼仪为操纵，\\
善意说真话却成冒犯，\\
完全读不懂空气
\end{cellvarwidth}\tabularnewline
\cmidrule{2-7}
 & \multirow{2}{*}{Demon} & 毁灭一切 & \multirow{2}{*}{ENTJ, ESTJ} & \begin{cellvarwidth}[t]
\centering
彻底崩溃\\
严重侵犯\\
极度绝望
\end{cellvarwidth} & \begin{cellvarwidth}[t]
\centering
名誉屠杀\\
Character Assassination
\end{cellvarwidth} & \begin{cellvarwidth}[t]
\centering
情感操纵与社会死亡，\\
情感爆发
\end{cellvarwidth}\tabularnewline
\cmidrule{3-3}\cmidrule{5-7}
 &  & 整合 &  & N/A & \begin{cellvarwidth}[t]
\centering
仁慈领袖\\
Benevolent Ruler
\end{cellvarwidth} & \begin{cellvarwidth}[t]
\centering
放下权力的傲慢，\\
展现出真正的人性光辉，\\
通过情感连接凝聚人心
\end{cellvarwidth}\tabularnewline
\bottomrule
\end{tabular}

}
\end{table}
